% Author-only replacement snippets. Not a standalone manuscript.
% Insert at the positions documented in MANUSCRIPT_EDITS.md.
% \nolinkurl is provided by url/hyperref; the manuscript already uses hyperlinks.
% Use the manuscript's existing labels for cross-references when available.

% Page 10: replace the existing Reproducibility Statement paragraph.
The appendix proof roadmap locates the complete mathematical proofs, and
Appendix~A states the imported certificate interfaces and prior results.
The supplement provides the scalar evaluator, reference Bellman recursions,
figure data, and reproduction commands.
\nolinkurl{NUMERICAL_REPRODUCIBILITY.md} distinguishes floating-point
computations from interval certificates. The phase-minimum and dyadic-corner
certificates are specified in Appendices~D.8 and~D.7, respectively.
The supplement also includes Lean~4 proofs of the scalar and combinatorial
results identified in \nolinkurl{lean/docs/THEOREMS.md}.
\nolinkurl{LEAN_FORMALIZATION.md} documents the pinned environment,
verification commands, and scope. The external smooth-convex GD certificates
and their tightness interpretation remain outside the formalization; the
level-17 decimal minimum enclosure is checked separately by interval arithmetic.

% Page 8: replace only the last sentence of the Figure 1 caption.
% Also re-include the corrected figures/psi_phase_compact.pdf and rebuild.
Numerical methods, reproduction commands, and the distinction between
finite-scale plots, interval certificates, and Lean verification are documented
in the supplementary \nolinkurl{NUMERICAL_REPRODUCIBILITY.md}.

% Page 11: optional sentence after the proof roadmap.
A statement-by-statement correspondence between the scalar and combinatorial
results and their Lean declarations is provided in
\nolinkurl{lean/docs/THEOREMS.md}; \nolinkurl{LEAN_FORMALIZATION.md}
documents the verification commands and scope limitations.

% Page 11: optional scope sentence in the Appendix A introduction.
The Lean development verifies the scalar and tree-based statements derived
from these interfaces. The smooth-convex function classes, gradient-descent
trajectories, and imported composition and tightness certificates recalled
in this appendix remain outside the formalization.

% Page 36: add after equation (65) and its certificate explanation.
Proposition~D.10 is formalized in Lean. The displayed level-17 decimal
enclosure is obtained separately by the interval-arithmetic procedure in
\nolinkurl{scripts/check_phase_certificate.py}; these particular decimal
endpoints are not part of the Lean certificate.
